\documentclass{article}
\usepackage[utf8]{inputenc}
\usepackage{amsmath, amsthm}
\usepackage{parskip}

\usepackage[french]{babel}

\usepackage{XCharter}
\usepackage[charter]{mathdesign}


\begin{document}
\title{MG tutorial notes}
\maketitle

\section{Relaxation methods}

\subsection*{Notations}

\begin{align*}
    \text{Problem:}          \quad & Au = f \\
    \text{Error:}            \quad & e = u - v \\
    \text{Residual:}         \quad & r = f - Av \\
    \text{Residual Problem:} \quad & Ae = r
\end{align*}

\subsection*{Relaxation Methods}

Decompose $A$ into $M - N$, where $M$ is easy to invert.

\[
    Au = f \implies Mu = Nu + f
\]

$u$ is a stationary point:
\[
    u = M^{-1}Nu + M^{-1}f
\]

With $R = M^{-1}N$, we use the iteration:
\[
    v^{(k+1)} = Rv^{(k)} + M^{-1}f
\]

Using $M^{-1}f = u - Ru$, we get:
\[
    e^{(k+1)} = Re^{(k)} = R^{k}e^{(0)}
\]

If $\|R\| < 1$, the error is forced to zero.

\subsubsection*{Examples (for matrices with nonzero diagonal)}
\begin{itemize}
    \item \textbf{Jacobi:} $A = D - L - U \quad M = D$
    
    \[
        Dv^{(k+1)} = (L + U)v^{(k)} + f
    \]
    \[
        a_{ii}v^{(k+1)}_{i} = \sum_{j \neq i}a_{ij}v^{(k)}_{j} + f_{i}
    \]
    solve for $v_{i}$ using the previous estimate

    \item \textbf{Gauss-Seidel:} $M = D - L$
    \[
        (D-L)v^{(k+1)} = Uv^{(k)} + f
    \]
    
    solve for $v_{j}$ using the previous estimate and already computed values
\end{itemize}

\subsubsection*{Issue \& strength}
Very slow convergence rate: good at the beginning, then stagnates.
Able to remove high-frequency/oscillatory error, but can't deal with low frequency/smooth one.
% \[
%     e \approx Re \implies (I - M^{-1}N)e \approx 0 \implies Ae \approx 0
% \]
% The error is in the kernel of $A$.
\[
    e^{(k+1)} = R^{k}e^{(0)}
\]

\begin{align*}
    R & = D^{-1}(L+U) \\
    & = D^{-1}(D - A) \\
    R & = I - \frac{1}{2}A
\end{align*}

In a 1D poisson with $n-1$ points, the eigenvalues of $A$ are $\lambda_{k} = 4\sin^{2}(\frac{k\pi}{2n})$

\subsubsection*{Coarse grid}

Starting point: If you subsample a signal, you increase the frequency of its modes.

idea for a recursive procedure:

- Relaxation until the error is smooth

- project the error on a smaller space where it becomes oscillatory again

- Relaxation \dots

Need a way to link the grids between eachother

\begin{itemize}
    \item \textbf{Interpolation/Prolongation:} coarse -> fine
    \item \textbf{Projection/Restriction:} fine -> coarse
\end{itemize}

Linear interpolation works well when the signal is smooth

what problem do we address on the coarse grid?
\[
A_{c}e_{c} = r_{c}
\]

In the geometric case, the physics dictate the operator: you can compute $A_{c}$ on the coarse grid

In the algebraic case, assuming the interpolation
\begin{align*}
    e \approx Pe_{c} \implies Ae = r & \approx APe_{c} \\ 
    P^{T}r & \approx P^{T}APe_{c}
\end{align*}
\begin{align*}
    A_{c}e_{c} & = r_{c} \\ 
    A_{c} & = P^{T}AP
\end{align*}


\section{Basic multigrid}



\end{document}